% !TEX root = ../TCC - CD e IA.tex
\chapter{Objetivos}\label{Objetivos e Organização do Trabalho}

A partir da problemática identificada \textemdash{} a distância entre o grande volume de dados de engajamento gerados nas mídias sociais e a carência de metodologias que os convertam em decisões práticas de comunicação \textemdash{} este trabalho estabelece um conjunto de objetivos que orientam o desenvolvimento e a validação de uma ferramenta de apoio à decisão para assessorias políticas. Diferentemente de abordagens que se encerram na descrição da opinião pública, aqui as técnicas de Processamento de Linguagem Natural (PLN) e de clusterização são o \emph{motor} de um instrumento orientado à ação: um funil de engajamento que permite diagnosticar a saúde da audiência de um perfil e priorizar o que produzir, o que amplificar e o que abandonar. A análise não se restringe aos comentários dos usuários; abrange as três fontes textuais disponíveis nas publicações \textemdash{} os \textbf{comentários}, as \textbf{legendas} (\emph{captions}) e as \textbf{transcrições} da fala dos vídeos \textemdash{} articuladas às \textbf{métricas de engajamento}, de modo a contrastar o discurso produzido pela assessoria (legendas e transcrições) com a reação do público (comentários). Os objetivos delineiam as etapas desse percurso, da estruturação do pipeline de dados até a tradução dos resultados em recomendações acionáveis, articulando os construtos de \emph{Consumers' Online Brand-Related Activities} (COBRAs) \cite{MUNTINGA2011} e o funil de \emph{marketing} de crescimento \cite{FEROZ2024} ao estudo de caso dos governadores do Brasil.

\section{Objetivo geral}

Desenvolver e validar uma metodologia de modelagem híbrida baseada em Processamento de Linguagem Natural (PLN) \textemdash{} combinando análise de sentimentos, modelagem de tópicos e clusterização sobre as diferentes fontes textuais e métricas das publicações \textemdash{} e traduzi-la em uma ferramenta de apoio à decisão que oriente estratégias de crescimento (\emph{growth}) de assessorias políticas na análise de publicações e comentários de redes sociais.

\section{Objetivos específicos}

Este trabalho tem os seguintes objetivos específicos:

\begin{enumerate}
    \item Estruturar um pipeline de processamento de dados textuais e quantitativos extraídos do Instagram, contemplando a coleta, limpeza, normalização e preparação das três fontes textuais das publicações \textemdash{} comentários, legendas e transcrições dos vídeos \textemdash{} bem como das respectivas métricas de engajamento.

    \item Aplicar técnicas de análise de sentimentos para identificar as polaridades expressas, distinguindo o sentimento dos comentários dos usuários \textemdash{} tanto em publicações do \emph{feed} quanto em \emph{Reels} \textemdash{} do sentimento veiculado pelas legendas e transcrições produzidas pelas assessorias.

    \item Implementar a modelagem de tópicos para extrair os principais temas em discussão, contrastando os temas emergentes nos comentários dos usuários com os temas do discurso oficial presente nas legendas e transcrições dos vídeos.

    \item Aplicar algoritmos de clusterização em diferentes granularidades, agrupando por semelhança tanto as publicações \textemdash{} \emph{Reels} e posts do \emph{feed} \textemdash{} quanto os perfis dos governadores, de modo a revelar padrões de desempenho e de comportamento comunicacional.

    \item Validar a eficácia da metodologia proposta por meio de um estudo de caso real, aplicando-a de forma integrada aos perfis, às publicações e às interações dos governadores do Brasil.

    \item Comparar os resultados da abordagem híbrida com os de abordagens tradicionais ou isoladas, evidenciando o ganho analítico obtido ao integrar sentimento, tópicos e clusterização sobre as múltiplas fontes textuais e níveis de agregação.

    \item Mapear os resultados da modelagem em um funil de engajamento acionável, articulando os níveis de COBRAs \cite{MUNTINGA2011} aos estágios do funil de crescimento \cite{FEROZ2024} e cruzando o discurso das assessorias (legendas e transcrições) com a resposta do público (comentários e métricas).

    \item Derivar, a partir desse mapeamento, recomendações práticas de \emph{growth} por perfil e por tipo de publicação, evidenciando como a ferramenta orienta a tomada de decisão de uma assessoria política.
\end{enumerate}





